\chapter{Conclusion}
\label{ch:conclusion}

This project built a teleoperation system for two seven-degree-of-freedom
arms mounted on a worn backpack frame, commanded by a second person through
an instrumented mannequin master or VR controllers, with a safety
architecture appropriate to a manipulator standing beside somebody who is not
commanding it and a shared-autonomy layer above it -- and ran a first pilot
comparison of direct and assisted control with real operators.

\section{Against the objectives}

A master was built and characterised: six of fourteen channels alive, the
repair identified to two specific joints. The command mapping degrades
predictably -- vertical and fore-aft track correctly, lateral does not, and
no recalibration fixes it. The safety architecture handled fourteen of
fourteen injected faults with every blocker observable by name. The
workspace characterisation found the two arms share no reachable point
under the as-built mount, identified a buildable mount correction, and that
correction was applied and re-verified. Ten-repeat, known-answer
verification ran throughout, catching twenty-six measurement defects before
they entered this report as findings. And a first pilot ran under ethical
approval with four operators across three tasks, showing shorter completion
times and fewer re-grips under assistance with no measured cost in hand
travel or responsiveness.

\section{What was not achieved}

No result here was validated on a physical Kinova arm carrying a load. The
full two-person study -- operator and wearer, measured together under the
same autonomy -- was not run; the pilot substitutes a smaller, operator-only
comparison, made possible once VR controllers offered a route around the
master's damaged channels. The planning report's haptic objective was not
delivered, and is not deliverable through this platform's control interface
without a hardware change.

\section{Further work}

In priority order: repairing the two identified master channels, which would
make the mannequin master usable as the direct-control baseline the full
study needs; calibrating the wrist-camera-to-arm transform, the largest
unmeasured term in the positioning budget and the reason no open-loop grasp
has yet closed on hardware; correcting the clearance check to sample link
geometry rather than link origins, the one item on this list that is a
safety defect rather than an incompleteness; and running the full two-person
protocol of \cref{app:study} with a separate wearer, now that ethical
approval and a working input path both exist, to measure what this project
was ultimately built to ask -- the exchange rate between what shared
autonomy buys an operator and what it costs the person carrying the arms.

\section{Closing}

The most transferable outcome of this work is not the mapping, the safety
architecture, or the pilot's result. It is the discipline that produced them:
twenty-six times a result that looked like a system failure turned out to be
a defect in the measuring tool, and the habit of establishing which one is in
hand before acting on it is the part of this project most worth carrying
into the study it was built to enable.
